\section{Formulario}
    \subsection{Trigonometria}\label{Trigonometria}
        \begin{enumerate}
            \item {
                $\sin^2(\alpha) + \cos^2(\alpha) = 1$
            }
            \item {
                $\cos(\alpha)=\pm\frac{1}{\sqrt{1+\tan^2(\alpha)}}$
            }
            \item {
                $\sin(\alpha)=\pm\frac{\tan(\alpha)}{\sqrt{1+\tan^2(\alpha)}}$
            }
            \item {
                $sinc(\alpha)\triangleq\frac{\sin(\pi\alpha)}{\pi\alpha}$
                É un $\sin(\alpha)$ smorzato secondo $\frac{1}{x}$ che si annulla in \\ $k\pi: k\in\mathbb{Z}$
                \begin{figure}[htp]
                    \centering
                    \includegraphics[width=4cm]{media/uwu.png}
                    \caption{grafico $sinc(\alpha)$}
                    \label{fig:grafico sinc}
                \end{figure}
            }
        \end{enumerate}
        \subsubsection{Formule di addizione}\label{Trigonometria_Addizione}
            \begin{enumerate}
                \item {
                    $\cos(\alpha \pm \beta) = \cos(\alpha)\cos(\beta) \mp \sin(\alpha)\sin(\beta)$
                }
                \item {
                    $\sin(\alpha \pm \beta) = \sin(\alpha)\cos(\beta) \pm \sin(\beta)\cos(\alpha)$
                }
                \item {
                    $\tan(\alpha \pm \beta) = \frac{\tan(\alpha) \pm \tan(\beta)}{1 \mp \tan(\alpha)\tan(\beta)} $
                }
            \end{enumerate}
        \subsubsection{Formule di duplicazione}\label{Trigonometria_Duplicazione}
            \begin{enumerate}
                \item {
                    $\sin(2\alpha) =2\sin(\alpha)\cos(\alpha)$ 
                }
                \item {
                    $
                        \cos(2\alpha)
                        \begin{cases}
                            \cos^2(\alpha) - \sin^2(\alpha) \\
                            2\cos^2(\alpha)-1\\
                            1-2\sin^2(\alpha)
                        \end{cases}
                    $
                }
                \item {
                    $\tan(2\alpha) =\frac{2\tan(\alpha)}{1-\tan^2(\alpha)}$ 
                }
            \end{enumerate}
            \subsubsection{Formule di bisezione}\label{Trigonometria_Bisezione}
                \begin{enumerate}
                    \item {
                        $\sin(\frac{\alpha}{2}) =\pm\sqrt{\frac{1-\cos(\alpha)}{2}}$ 
                    }
                    \item {
                        $\cos(\frac{\alpha}{2}) =\pm\sqrt{\frac{1+\cos(\alpha)}{2}}$ 
                    }
                    \item {
                        $
                            \tan(\frac{\alpha}{2})
                            \begin{cases}
                                \sqrt{\frac{1-\cos(\alpha)}{1+\cos(\alpha)}} \\
                                \frac{1-\cos(\alpha)}{\sin(\alpha)}\\
                                \frac{\sin(\alpha)}{1+\cos(\alpha)}
                            \end{cases}
                        $
                    }
                \end{enumerate}
    \subsection{Segnali Comuni}\label{Segnali Comuni}
